\begin{textsample}{NEG Dim 2 – mg – Score -1.00 – mg/mg\_em\_43.txt}  \label{ex:f2_neg_100}
3.1.3.1.1 Leitura \textbf{literária} e construção de sujeitos críticos A BNCC traz habilidades a serem desenvolvidas no Campo Artístico \textbf{Literário}, as quais buscam contribuir para a formação integral de estudantes, proporcionando -lhes vivências estéticas por meio da interação com elementos simbólico-culturais e com recursos expressivos constituintes da linguagem \textbf{literária}.
Por meio dessa interação, o estudante amplia o repertório linguístico-cultural e a forma como percebe os sonhos, as realidades e os revezes alheios, reconhecendo a si mesmo, o outro e o nós, (re)construindo o sentido do texto e de sua realidade.
Assim, desenvolve a capacidade de análise, constituindo -se como um sujeito leitor e protagonista.
Desse modo, a literatura, em quaisquer de suas funções – lúdica, pragmática, catártica ou estética – busca a construção de um ser potencialmente reflexivo, ético, participativo e atuante em diversos contextos sociais, permeados ou não pela cultura digital.
Na BNCC, a literatura busca contribuir para que as juventudes estabeleçam posturas críticas em relação aos diversos textos - orais ( nos saraus, nos slams, nos duelos de MCs, nos festivais \textbf{literários}, em rodas de conversa, nos podcasts ), escritos ( fanzines, e-zines, memes, blogs, resenhas, playlists \textbf{comentadas}, livros impressos ) criados por outros ou por eles mesmos, levando em consideração elementos tais como contexto de produção, autoria, intertextualidades possíveis, momentos culturais vigentes quando da produção e, ainda, a intencionalidade dessa produção \textbf{literária}.
Longe de negar a existência de ideologias, o campo \textbf{literário} na BNCC visa à formação de leitores capazes de identificá -las, de compreendê -las e de confrontá -las, constituindo -se como um sujeito de multiplicidade de culturas, de textos, de contextos ; apropriando -se de novos elementos éticos, estéticos, culturais e simbólicos.
Capazes de atuar crítica e eticamente em sua realidade.

% matched lemmas: comentar, literário
\end{textsample}
